\chapter{Results and Evaluation}
\label{ch:results}

At the mid-term point, FoodLink has not been deployed or used by real
donors, organisations or couriers. Its evaluation therefore rests on
\emph{verification}: checking that the implemented behaviour matches the
requirements, and that the security, privacy and integrity properties hold
under adversarial and concurrent use. This chapter describes the testing
strategy, reports the results actually obtained while preparing this report,
and analyses what they show, including the system's current limitations.

\section{Testing Methodology}
\label{sec:testing}

Testing follows the project's principle that a property which matters must be
held by an automated test~\cite{pressman2015}. Every fix from the audits
(Section~\ref{sec:analysis}) arrived with regression tests that fail against
the code before the fix. The strategy has four parts:

\begin{itemize}
  \item \textbf{Unit Testing:} The ranking arithmetic is tested directly,
    without HTTP or a database, in \code{test_matching_scores.py} and
    \code{test_requirement_matching.py}. These cover each criterion's formula,
    the eligibility gates, unit handling, the blurring of positions and the
    requirement tie-break. Configuration validation and the sliding-window rate
    limiter are also unit-tested (\code{test_config.py},
    \code{test_rate_limit.py}). In the frontend, Vitest unit tests cover the
    pure library modules: wire-to-domain adapters, time and urgency, distance
    and location display, per-account impact, the HTTP client, and the photo
    resizing pipeline, with its browser image functions injected because jsdom
    has no canvas.
  \item \textbf{Integration Testing:} Most backend tests exercise the complete
    HTTP stack through FastAPI's test client against a fresh in-memory SQLite
    database per test, with no mocks (D-17). A test registers or seeds real
    accounts, signs in, and calls the API exactly as the frontend does. This
    is how authorization boundaries are tested: every role is tried against
    every scoped read and every transition. Concurrency tests use a
    file-backed database and two genuine transactions interleaved by hand, so
    that a lost-update race is actually reproduced and then shown to be
    refused. Migration tests build a database from the Alembic revision
    history and assert that it matches the models. Frontend component tests
    render screens with Testing Library in jsdom to check what each screen
    shows and, just as importantly, what it must not show.
  \item \textbf{Performance Testing:} No load, stress or response-time testing
    has been performed. Performance-related work so far consists of design
    measures (eager loading, bounded payloads, a 500-row list cap) and one
    browser measurement of the photo-resizing pipeline recorded during Task~40
    (Section~\ref{sec:module-portals}).
  \item \textbf{User Testing:} Not yet conducted. The proposal's pilot, with
    campus donors, recipient organisations and student
    volunteers~\cite{foodlinkProposal}, remains future work
    (Chapter~\ref{ch:feasibility}).
\end{itemize}

\paragraph{Continuous integration.} A GitHub Actions workflow
(\code{.github/workflows/ci.yml}) runs on every push to the main branch and on
every pull request. It has two parallel jobs. The \emph{backend} job, on
Python~3.13, runs \code{pytest code/tests}, then \code{alembic upgrade head}
and \code{alembic check} on a throwaway database to catch a model change
committed without a migration. The \emph{frontend} job, on Node.js~20, runs
\code{npm ci}, \code{npm test} and \code{npm run build}; the build step
includes the TypeScript compiler. The workflow deploys nothing and holds no
secrets (D-25).

Table~\ref{tab:testfiles} groups the 27 backend test files by the property
they protect.

\begin{table}[htbp]
\centering
\small
\begin{tabular}{|L{3.0cm}|L{10.3cm}|}
  \hline
  \textbf{Area} & \textbf{Backend test files} (in \code{code/tests/}) \\
  \hline
  End-to-end paths, accounts, configuration & \code{test_api.py}, \code{test_auth_admin.py}, \code{test_config.py}, \code{test_rate_limit.py}, \code{test_donation_rate_limit.py}, \code{test_migrations.py}, \code{test_recipient_reverification.py} \\
  \hline
  Input validation & \code{test_donation_input_bounds.py}, \code{test_patch_null_semantics.py} \\
  \hline
  Read and disclosure scopes & \code{test_donation_reads.py}, \code{test_recipient_reads.py}, \code{test_volunteer_reads.py}, \code{test_requirement_reads.py}, \code{test_courier_pickup_disclosure.py} \\
  \hline
  Lifecycle and concurrency & \code{test_lifecycle_authorization.py}, \code{test_pickup_release.py}, \code{test_cancellation_reliability.py}, \code{test_courier_claim.py}, \code{test_lifecycle_concurrency.py}, \code{test_available_donations_deadline.py}, \code{test_volunteer_delivery_history.py}, \code{test_requirement_lifecycle.py} \\
  \hline
  Matching and location privacy & \code{test_matching_scores.py}, \code{test_match_score_consistency.py}, \code{test_match_distance_privacy.py}, \code{test_donation_privacy_scope.py}, \code{test_requirement_matching.py} \\
  \hline
\end{tabular}
\caption{Backend test suite organised by the property under test}
\label{tab:testfiles}
\end{table}

\section{Performance Metrics}
\label{sec:metrics}

All results in Table~\ref{tab:validation} were obtained by running the checks
on the repository at commit \texttt{e588b35} on 15 September 2026 while
preparing this report. The backend database was redirected to a temporary
file so that no project file was changed.

\begin{table}[htbp]
\centering
\small
\begin{tabular}{|L{5.0cm}|L{3.1cm}|L{5.2cm}|}
  \hline
  \textbf{Metric} & \textbf{Target} & \textbf{Achieved} \\
  \hline
  Backend automated tests (\code{pytest code/tests}) & All pass & 495 passed, 0 failed, 1 library deprecation warning; 337\,s \\
  \hline
  Frontend automated tests (\code{npm test}) & All pass & 156 passed, 0 failed, in 18 files \\
  \hline
  Frontend type check (\code{tsc --noEmit}) & No errors & No errors \\
  \hline
  Migration drift (\code{alembic upgrade head}, \code{alembic check}) & No drift & Upgraded to \code{ae4636b1e6d4}; no new operations detected \\
  \hline
  Lifecycle states enforced & 9 & 9 \\
  \hline
  API routes implemented & All proposal workflows & 27 routes across 5 routers plus health \\
  \hline
  Role portals & 4 & 4 web portals, plus phone-layout screens \\
  \hline
  P1 findings of the 10 September 2026 audit & All fixed & 5 of 5 fixed, each with regression tests \\
  \hline
\end{tabular}
\caption{Validation results at mid-term}
\label{tab:validation}
\end{table}

The frontend production build (\code{npm run build}) was not re-run for this
report, because it writes build output into the repository. It runs in CI on
every push. No response-time or throughput figures are reported, because none
have been measured.

The proposal defined outcome metrics for a pilot~\cite{foodlinkProposal}.
Table~\ref{tab:proposal-metrics} records how far each has progressed. The
system can now \emph{compute} the ledger-based ones. None has been
\emph{measured}, because no real usage has taken place.

\begin{table}[htbp]
\centering
\small
\begin{tabular}{|L{3.4cm}|L{3.4cm}|L{6.5cm}|}
  \hline
  \textbf{Proposal metric} & \textbf{Proposal target} & \textbf{Mid-term status} \\
  \hline
  Median time-to-claim & $\leq$ 45 minutes & Instrumented in \code{GET /api/metrics}; not measured \\
  \hline
  Rescue rate & $\geq$ 80\% & Instrumented; not measured \\
  \hline
  Expiry-loss rate & Counter-metric & Instrumented; depends on the manual sweep having run \\
  \hline
  Handover time & Tracked & Instrumented (median acceptance to delivery); not measured \\
  \hline
  User clarity (1--5 rating) & Tracked & Not instrumented; no feedback feature exists \\
  \hline
  Service availability & $\geq$ 99\% in pilot & Not applicable yet; nothing is deployed \\
  \hline
\end{tabular}
\caption{Evaluation metrics from the proposal and their mid-term status}
\label{tab:proposal-metrics}
\end{table}

\section{Analysis}
\label{sec:analysis}

\begin{itemize}
  \item \textbf{Did the project meet its objectives?} Objectives~1 to~5 and~7
    are met and verified (Table~\ref{tab:objectives}). Objective~6 is met as
    instrumentation but cannot yet produce meaningful numbers without real use.
    Objective~8 (deployment and pilot) has not started.
  \item \textbf{Where did the work go beyond the proposal?} The proposal
    described a first iteration of posting, acceptance, courier status updates
    and a shared timeline, with ranking and requirements as later
    deliverables~\cite{foodlinkProposal}. Both later deliverables are already
    implemented. So is a set of properties the proposal did not name, which
    emerged from the audits: requirement-aware tie-breaking, per-reader
    disclosure scopes for donor and organisation locations, race-safe
    conditional writes for every transition, abuse limits, and reliability
    accounting that is not distorted by cancellations.
  \item \textbf{What challenges were encountered?} The hardest problems were
    not features but properties that only appear under adversarial or
    concurrent use. Each was found by an audit and reproduced against the
    running application before it was fixed. Table~\ref{tab:challenges}
    summarises the most significant.
  \item \textbf{How were they resolved?} Each fix was the smallest change that
    made the property hold everywhere rather than only on the reported path.
    It was recorded as a decision with its reasoning and protected by tests
    that fail without it. None of the five P1 fixes required a database schema
    change.
\end{itemize}

\begin{table}[htbp]
\centering
\footnotesize
\begin{tabular}{|L{4.0cm}|L{5.4cm}|L{3.9cm}|}
  \hline
  \textbf{Challenge (audit finding)} & \textbf{Resolution} & \textbf{Evidence} \\
  \hline
  Newly self-registered organisation and courier accounts could read every open donation's exact pickup pin, address and donor name (P1-1) & Open pool restricted to verified organisations (D-53); couriers see a coarse area until they claim (D-57) & \code{test_donation_reads.py}, \code{test_courier_pickup_disclosure.py} \\
  \hline
  Verification survived an organisation renaming or moving itself (P1-2) & A real change to name or coordinates clears verification in the same commit (D-54) & \code{test_recipient_reverification.py} \\
  \hline
  Two organisations accepting at the same moment both succeeded, even on SQLite (P1-3) & Every status write made conditional on the expected state; the loser receives 409 and its side effects roll back (D-55) & \code{test_lifecycle_concurrency.py} \\
  \hline
  Donor cancellations lowered the organisation's reliability; donors could cancel collected food (P1-4) & Cancellation withdraws the counted acceptance; donors cannot cancel after pickup (D-58) & \code{test_cancellation_reliability.py} \\
  \hline
  Ordinary phone photos made donation posting fail against the image cap (P1-5) & Browser-side resize and JPEG re-encode before upload (D-56) & \code{lib/__tests__/image.test.ts} \\
  \hline
  Ranking results revealed exact distances, from which organisation addresses could be triangulated (HA-3) & Organisation positions blurred to a 0.01$^\circ$ grid for other readers; distances and frozen scores withheld (D-45, D-47) & \code{test_match_distance_privacy.py} \\
  \hline
  Unlimited donation posting allowed spam and probing of the radius gate (P2-1) & Per-account and per-address limits (D-59) & \code{test_donation_rate_limit.py} \\
  \hline
  Two quantity criteria were mathematically the same criterion counted twice, and ignored units (HA-4, HA-5) & Capacity headroom redefined in absolute meals; non-meal units not assessed (D-42) & \code{test_matching_scores.py} \\
  \hline
\end{tabular}
\caption{Principal challenges found by audit and how they were resolved}
\label{tab:challenges}
\end{table}

\subsection{Current Limitations}
\label{sec:limitations}

The following limitations are verified against the repository and are stated
so that the system is not described as more complete than it is.

\begin{enumerate}
  \item \textbf{No deployment.} Nothing is hosted. There is no deployment
    configuration, and PostgreSQL, although supported in principle, is untested
    and has no driver pinned.
  \item \textbf{No notifications.} Participants learn about new donations and
    status changes only by opening the portal.
  \item \textbf{Missing web controls.} Cancellation, pickup release and account
    management work through the API but have no portal screen. Neither the
    registration form nor the organisation profile screen collects map
    coordinates. An organisation created through the portal therefore cannot be
    ranked, since coordinates are an eligibility gate, until its pin is set
    through \code{PATCH /api/recipients/me}. The same applies to organisations
    created by an administrator. In practice, only organisations with
    coordinates, such as the seeded demonstration organisations, appear in
    rankings today.
  \item \textbf{Manual expiry.} The expiry sweep must be triggered by an
    administrator. Its bulk write does not use the conditional guard that every
    other status write uses. An administrator's \code{ACCEPTED} $\to$
    \code{EXPIRED} leaves the acceptance counted.
  \item \textbf{Session security.} No token revocation, refresh token,
    multi-factor authentication or content security policy. Rate-limit counters
    are per process and reset on restart.
  \item \textbf{No maps.} Distance is great-circle, and travel time is a flat
    20\,km/h estimate; there is no road routing. The pickup pin is entered
    through browser geolocation or typed coordinates, with no map picker.
  \item \textbf{Heuristic weights.} The matching weights and the reliability
    prior are reasoned judgements, not values tuned on outcome data.
  \item \textbf{Inline photos.} Photos are stored in the donation row and
    returned with every list, and every write re-downloads up to 500 donations.
  \item \textbf{Scale.} Recipients are loaded and filtered in Python for each
    ranking. Metrics are computed in Python over every donation. SQLite allows
    a single writer.
  \item \textbf{Validation gaps.} Account and organisation profile text is
    unbounded. An unexpected server error returns an empty 500 response, which
    the frontend presents as a connection problem.
  \item \textbf{Demonstration residue.} The phone-layout camera screen
    simulates reading a photo and submits a fixed address, and one mobile
    profile screen shows a hard-coded email. The \code{/ngo/available/:id} deep
    link does not select the donation.
  \item \textbf{Metric caveats.} The completed-meals total adds quantities
    across different units. The expiry-loss rate is accurate only after the
    sweep has run.
\end{enumerate}
